% RESTORED VERSION
\documentclass[10pt,a4paper]{article}

\usepackage[T1]{fontenc}
\usepackage[utf8]{inputenc}
\usepackage[french]{babel}
\usepackage{geometry}
\geometry{top=0.8cm,bottom=2.5cm,left=1.8cm,right=1.8cm}
\usepackage{lmodern}
\usepackage{microtype}
\usepackage{parskip}
\setlength{\parskip}{6pt}
\usepackage{graphicx}
\usepackage{tabularx}
\usepackage{array}
\usepackage{booktabs}
\usepackage{fancyhdr}
\usepackage{hyperref}
\hypersetup{colorlinks=false,pdfborder={0 0 0}}
\usepackage{enumitem}
\usepackage{eurosym}
\usepackage{titlesec}

% ---------- section style ----------
\titleformat{\section}
  {\normalsize\bfseries}
  {}{0em}{}[\vspace{-2pt}\rule{\linewidth}{0.8pt}\vspace{2pt}]
\titlespacing{\section}{0pt}{12pt}{6pt}
\titleformat{\subsection}
  {\normalsize\bfseries\itshape}{}{0em}{}
\titlespacing{\subsection}{0pt}{8pt}{3pt}

% ---------- footer ----------
\pagestyle{fancy}
\fancyhf{}
\renewcommand{\headrulewidth}{0pt}

\begin{document}
\thispagestyle{fancy}

% ================================================================
%  TITLE
% ================================================================
\begin{center}
  {\Large\bfseries Présentation de l'entreprise et de l'équipe SAA}\\[6pt]
  {\normalsize\itshape Dossier de candidature CIFRE --- Yassine MANNAI}
\end{center}

\vspace{0.4cm}

% ================================================================
\section{Présentation de Generali France}
% ================================================================

Generali France est l'un des principaux assureurs de l'Hexagone, s'appuyant
sur l'esprit pionnier d'un groupe international fondé en 1831.  L'entreprise
opère sur l'ensemble des segments de l'assurance (Vie, Dommages, Santé,
Prévoyance et Retraite) avec une ambition de \og{}Lifetime Partner\fg{} pour
ses clients.

% ----------------------------------------------------------------
\subsection{Indicateurs clés et volume d'activité (2025)}
% ----------------------------------------------------------------
\begin{itemize}[leftmargin=1.6em,itemsep=1pt]
  \item \textbf{Chiffre d'affaires :} environ 20 milliards d'euros.
  \item \textbf{Résultat opérationnel :} environ 1~400 millions d'euros (normes IFRS~17/9).
  \item \textbf{Actifs sous gestion :} environ 124 milliards d'euros.
  \item \textbf{Effectif total :} 9~300 collaborateurs et salariés commerciaux.
  \item \textbf{Rayonnement :} 8 millions de clients et un réseau de près de 800 agents généraux.
\end{itemize}

% ----------------------------------------------------------------
\subsection{Organisation et direction de l'encadrement}
% ----------------------------------------------------------------

L'encadrement industriel de la thèse est rattaché à la \textbf{Direction des
Investissements} de Generali France, placée sous la responsabilité de
\textbf{Cédrik d'Aviau de Ternay} (Directeur des Investissements).  Cette
direction pilote l'allocation stratégique du bilan de Generali France dans le
cadre de Solvabilité~II, en s'appuyant notamment sur l'outil propriétaire du
Groupe, \textbf{GPS} (\textit{Group Portfolio Solutions}).
Elle est organisée en deux équipes complémentaires : l'équipe \textbf{SAA}
(Allocation Stratégique d'Actifs), dédiée à la recherche quantitative et à
l'optimisation du bilan, et l'équipe \textbf{Pilotage}, chargée du suivi
opérationnel des allocations et de la coordination avec Generali Asset
Management (GenAM).

% ================================================================
\section{Équipe d'accueil : Allocation Stratégique d'Actifs (SAA)}
% ================================================================

Le candidat sera intégré au sein de l'\textbf{équipe Allocation Stratégique
d'Actifs (SAA)}, entité de recherche quantitative et de pilotage du bilan au
sein de la Direction des Investissements.  L'équipe regroupe des expertises en
finance quantitative, gestion actif-passif (ALM), modélisation actuarielle et
optimisation de portefeuille sous contraintes réglementaires.

\begin{itemize}[leftmargin=1.6em,itemsep=1pt]
  \item \textbf{Directeur des Investissements :} Cédrik d'Aviau de Ternay.
  \item \textbf{Responsable scientifique (Encadrant industriel) :}
        Mamy Ramamonjisoa, membre de l'équipe SAA.
  \item \textbf{Membres de l'équipe SAA :}
        Teddy Vaurs, Mohamed Koudiraty, Cesare Mele ;
        Raphaël Konan (consultant associé).
  \item \textbf{Doctorant CIFRE (candidat) :} Yassine Mannai.
\end{itemize}

% ================================================================
\section{L'équipe SAA comme service de recherche quantitative}
% ================================================================

L'équipe SAA constitue le pôle de recherche quantitative appliquée de la
Direction des Investissements.  Elle est chargée de :

\begin{itemize}[leftmargin=1.6em,itemsep=1pt]
  \item \textbf{Pilotage stratégique du bilan :} construction et révision de
        l'allocation stratégique d'actifs dans le cadre de Solvabilité~II, via
        l'outil propriétaire GPS (\textit{Group Portfolio Solutions}).
  \item \textbf{Stratégies de couverture :} étude et mise en \oe{}uvre des
        couvertures de portefeuille en scénario de stress, en lien avec les
        brokers et Generali Asset Management (GenAM).
  \item \textbf{Recherche et innovation :} intégration des approches de la
        finance quantitative contemporaine (modèles génératifs, proxy ML)
        dans les processus d'allocation, en vue notamment d'embarquer des
        modèles légers et efficaces directement dans la boucle d'optimisation.
  \item \textbf{Coordination multi-équipes :}
  \begin{itemize}[leftmargin=1.4em,itemsep=0pt]
    \item \textit{Équipe Pilotage} --- définition des bornes d'allocation par
          classe d'actifs et suivi de l'exécution via GenAM ;
    \item \textit{Direction des Risques} --- transmission des besoins de calcul
          SCR et intégration des résultats dans les optimisations ;
    \item \textit{Équipe Plan \& Reporting} --- accès aux bases comptables pour
          la calibration des modèles ;
    \item \textit{Équipe ALM} --- projections actif-passif nécessaires aux
          simulations de scénarios à long terme.
  \end{itemize}
\end{itemize}

% ================================================================
\section{Projet de thèse CIFRE}
% ================================================================

% ----------------------------------------------------------------
\subsection{Titre}
% ----------------------------------------------------------------

\begin{center}
  \itshape
  Modèles génératifs conditionnels aux régimes, proxys de capital de Volterra
  et dynamiques de rachat en champ moyen pour l'allocation d'actifs en assurance
\end{center}

% ----------------------------------------------------------------
\subsection{Contexte et motivation industrielle}
% ----------------------------------------------------------------

L'épisode de 2022 --- inversion simultanée de la corrélation
actions/obligations sur l'ensemble des marchés --- a révélé quatre
limitations structurelles des outils SAA actuels :
(i)~corrélations aveugles aux ruptures de régime ;
(ii)~calcul du SCR (Direction des Risques, LSMC, plusieurs heures)
incompatible avec des optimisations dynamiques et scalables sur
multiples scénarios ;
(iii)~actifs privés à valorisation lissée biaisant le risque réel ;
(iv)~absence de modèle de contagion pour les rachats corrélés.
La thèse vise à rendre opérationnel le \textbf{Total Portfolio Approach (TPA)}
pour un assureur sous Solvabilité~II.

% ----------------------------------------------------------------
\subsection{Axes de recherche}
% ----------------------------------------------------------------

\begin{enumerate}[leftmargin=2em,itemsep=4pt]
  \item \textbf{Simulation génératrice conditionnelle aux régimes (Axe~1).}
        Développement et comparaison de modèles génératifs
        (GAN conditionnel, VAE conditionnel, Score-Based Generative Models)
        pour simuler la distribution jointe des facteurs TPA sous chaque régime
        macroéconomique identifié en temps réel par un filtre séquentiel.
        Validation formelle de la structure de dépendance apprise par des tests
        de spécification sur copules conditionnelles (Fermanian, 2013, 2022)
        et établissement de garanties de convergence Wasserstein-2.

  \item \textbf{Proxy de SCR supervisé pour la décision dynamique (Axe~2).}
        Construction d'un \textit{tiny model} différentiable du capital
        réglementaire (SCR), conçu pour s'embarquer directement dans la boucle
        d'optimisation de l'équipe SAA --- sans dépendre des recalculs complets
        de la Direction des Risques (Prophet/MoSes, plusieurs heures).
        Exprimé en fonction des facteurs TPA, du régime courant et des états
        des passifs (dont les lags de NAV d'actifs privés), il offre une
        inférence inférieure à 10~ms avec des bornes d'erreur contrôlées ---
        condition nécessaire pour scaler vers des optimisations conjointes
        dynamiques sur de multiples scénarios de marché simultanés.

  \item \textbf{Rachats en champ moyen et allocation de capital (Axe~3).}
        Modélisation de la dynamique de rachat des assurés comme un équilibre
        de jeu à champ moyen (MFG) sur un réseau de contrats hétérogènes,
        conditionné au régime macroéconomique identifié par l'Axe~1.
        Le processus agrégé de rachat alimente directement le proxy SCR de l'Axe~2
        et permet de résoudre un programme d'allocation de capital sous contrainte
        Solvabilité~II explicite.
\end{enumerate}

% ================================================================
\section{Intérêt des travaux pour le développement de l'entreprise}
% ================================================================

La thèse s'inscrit dans le pilier \og{}Excellence\fg{} du plan stratégique
\textbf{Boost 2027}.  Pour Generali France, l'intérêt est triple :

\begin{enumerate}[leftmargin=2em,itemsep=4pt]
  \item \textbf{Résilience de l'allocation stratégique.}
        Les travaux fourniront un outil de simulation capable d'anticiper les
        ruptures de régime et d'adapter l'allocation du bilan avant les phases
        de stress, réduisant l'exposition aux pertes extrêmes dues à des
        dépendances non linéaires entre facteurs.

  \item \textbf{Accélération du calcul du capital réglementaire.}
        Le proxy SCR supervisé permettra d'intégrer le coût en capital
        directement dans les boucles d'optimisation en temps quasi-réel,
        rendant possible une SAA dynamique --- là où le recalcul complet
        (Prophet/MoSes) nécessite aujourd'hui plusieurs heures.

  \item \textbf{Gestion du risque de rachat.}
        Le modèle MFG fournit un processus de rachat agrégé calibré sur
        l'épisode de 2022 et intégré comme contrainte explicite dans le
        programme d'allocation de capital sous Solvabilité~II, remplaçant
        les tables de rachat déterministes actuelles.
\end{enumerate}

\medskip
\noindent Les trois modules sont conçus pour s'intégrer à \textbf{GPS}
(\textit{Group Portfolio Solutions}), offrant à Generali France un avantage
concurrentiel durable dans la gestion actif-passif et le pilotage du bilan.

\end{document}
